导数回答``此刻变化多快''，积分回答``这些变化攒到最后是多少''。两个问题看上去朝相反的方向问，微积分基本定理却说它们互为逆运算——正因为如此，求曲边下的面积才被降格成了``找一个导数等于被积函数的函数''。这条互逆关系是本章的主轴。

先看一眼它有多省事。速度 \(v(t)=2t\) 在 \([0,3]\) 上跑出多少位移？按定义，把 \([0,3]\) 切成 6 段，每段宽 \(0.5\)，用右端点速度 \(1,2,3,4,5,6\) 近似，得到 \((1+2+3+4+5+6)\times0.5=10.5\)；继续加密，这串数会稳定到 9。但如果你注意到 \(t^2\) 的导数正是 \(2t\)，直接写 \(3^2-0^2=9\) 就完了。整章要说清楚的，就是第二条路凭什么合法，以及它在什么地方不再合法。

对做数据的人来说，积分的现实分量比几何面积重得多。连续分布下的一切都是积分：概率是密度曲线下的面积，期望是加权积分，配分函数是归一化用的那个积分，累积分布函数就是密度的变上限积分。换元法带来的那个``密度因子''，正是后面概率里变量替换、多重积分的 Jacobian 行列式、以及 normalizing flow 里 \(\log\abs{\det J}\) 项的同一件事。反常积分是否收敛，则决定了一个重尾分布的期望到底存不存在。

\subsection{怎么读这一章}

第一节把定积分立起来：切细、取样、求和、取极限，以及这个极限凭什么与切法无关。第二节是全章的支点，把累积与变化率焊在一起。第三、四节讲技巧，但读它们的时候请盯住结构变换本身——换元在做坐标伸缩，分部在搬运导数，部分分式在把有理函数还原成已知的几块——具体题型不值得记。第五节收口，把同一套切片动作铺到面积、体积、概率与期望上。

若你只想拿走一件东西，拿第二节；若你为概率或机器学习而来，第三节的密度因子和第五节的密度与累积量最直接。

\subsection{本章篇目}

\begin{itemize}
\tightlist
\item \hyperref[sec:02-Calculus/04-Integration/01]{``从小块求和到定积分''}：为什么``图像围出一块面积''不算定义，可积性是极限的唯一性；
\item \hyperref[sec:02-Calculus/04-Integration/02]{``微积分基本定理：变化与累积互相撤销''}：累积函数的导数为什么就是被积函数，连续性用在哪一步；
\item \hyperref[sec:02-Calculus/04-Integration/03]{``换元与分部积分：反向识别求导规则''}：换元的实质是坐标伸缩因子，分部的实质是拿一个边界项换一个更好算的积分；
\item \hyperref[sec:02-Calculus/04-Integration/04]{``部分分式与反常积分''}：代数还原路径，以及收敛与否只取决于尾部衰减多快；
\item \hyperref[sec:02-Calculus/04-Integration/05]{``面积、体积、概率与累积量''}：同一个切片动作在几何与概率里的两副面孔。
\end{itemize}
